\chapter{总结与后续工作}

\section{工作总结}

WaterOS目前已经形成一条可运行实际用户程序的内核主路径。RISC-V64与LoongArch64分别
完成启动、异常、分页和设备接入，上层共用任务、内存、VFS/FS、IPC、网络及系统调用组件。
用户程序可以从ext4根文件系统装载，创建进程和线程，通过fd访问文件、终端与socket，并
在多核调度下运行。

项目开发中最费工作量的部分并不是增加接口数量，而是处理对象跨越组件之后的生命周期。
fork同时涉及任务、地址空间、fd、凭证和信号；退出需要让调度器停止使用任务，再由父进程
等待并回收资源；共享文件映射又把VMA销毁与页缓存写回联系起来。围绕这些路径，项目建立
了先初始化后发布、阻塞前登记等待、按引用关系销毁对象以及写回版本检查等约定。它们并非
所有问题的终点，但使故障能够落到具体对象和状态转换上，而不是依赖不断增加临时日志。

双架构实现也帮助我们划清了可复用机制与硬件事实之间的界线。平台接口可以统一“发送
IPI”或“切换地址空间”的调用形式，但不能抹平页表格式、固件服务和TLB维护能力的差别。
因此本文保留了LoongArch64远程TLB失效、完整权限模型、设备事件唤醒和文件系统异常恢复等
尚未完成的边界，没有将部分测例通过表述成完整Linux兼容。

\section{后续工作}

后续工作将优先处理会影响长期运行的基础问题，而不是继续扩大接口清单：

\begin{enumerate}
  \item 完成LoongArch64跨核TLB失效，并为地址空间进入、退出、解除映射和销毁建立双架构
  一致的并发测试。
  \item 继续检查高频fork后的内核堆、任务注册表、页表和fd侧表回收，增加按对象类型统计
  的高水位与回落指标，使泄漏能够在OOM之前被发现。
  \item 补充页缓存写回的并发测试，覆盖共享映射、truncate、rename、进程退出和写回失败，
  明确哪些错误可以返回给用户，哪些只能在异步回收阶段记录。
  \item 将socket等待从定时扫描逐步改为设备事件驱动，并区分协议栈推进、可读写状态变化
  与用户等待队列唤醒，减少空轮询。
  \item 完善uid/gid、inode访问控制、capability和TTY作业控制；对仍未实现的Linux语义
  继续显式返回错误，避免“接口存在但行为错误”。
  \item 把双架构构建、完整负载、fork压力和镜像一致性检查纳入固定回归流程，并让结果与
  commit、镜像摘要和工具版本绑定。
\end{enumerate}

\section{第三方代码说明}

WaterOS的顶层组合、组件接口、平台适配、任务与内存管理集成、VFS桥接、generic64系统
调用分发和测试编排由项目组维护。项目使用第三方crate完成协议栈、设备访问、设备树解析
和部分文件系统工作；部分上游代码以源码形式引入，并维护了适配WaterOS的修改。表
\ref{tab:third-party}列出与本文实现关系较大的依赖，精确版本以构建锁定结果为准。

\begin{longtable}{L{3.0cm}L{6.6cm}L{3.6cm}}
  \caption{主要第三方依赖}\label{tab:third-party}\\
  \toprule
  依赖 & 用途 & 使用方式 \\
  \midrule
  \endfirsthead
  \toprule
  依赖 & 用途 & 使用方式 \\
  \midrule
  \endhead
  another\_ext4 & 默认ext4读写后端；WaterOS维护了与本项目接口和写路径相关的修改。 &
  源码维护版本 \\
  ext4\_rs、ext4plus & ext4相关实现与备用接入路径。 & 源码引入 \\
  smoltcp 0.12.0 & IPv4 TCP/UDP协议栈基础。 & crates.io锁定版本 \\
  virtio-drivers 0.12.0 & VirtIO设备访问。 & crates.io锁定版本 \\
  fdt 0.1.5 & RISC-V启动时的设备树解析。 & crates.io锁定版本 \\
  sbi-rt 0.0.3 & RISC-V OpenSBI运行时调用。 & crates.io锁定版本 \\
  riscv 0.16.0 & RISC-V CSR与架构辅助。 & crates.io锁定版本 \\
  rlsf 0.2.2 & 内核堆分配器后端之一。 & crates.io锁定版本 \\
  \bottomrule
\end{longtable}

用户空间中的BusyBox、Nano-X、OpenJDK及其他软件有各自的上游来源和许可证，并由用户
根文件系统构建组件统一组织。这些程序用于验证内核行为，不计入内核自研实现。

\section{AI工具使用说明}

开发过程中使用了Claude Code、Codex和Cursor等AI辅助工具。AI主要参与功能代码生成、重复性
接口适配、跨组件调用关系检索、故障定位、定向测试编写以及开发文档整理。系统的组件边界、
接口取舍、兼容范围和最终合入决定由项目成员负责；AI生成或修改的实现仍需经过源码审查、
目标架构编译以及与修改范围相对应的运行测试，不能复现的内容不作为本文结论。

\begin{longtable}{L{2.8cm}L{2.3cm}L{9.1cm}}
  \caption{项目组件中的AI工具使用情况}\label{tab:ai-usage-by-module}\\
  \toprule
  组件 & 组件功能 & AI使用情况 \\
  \midrule
  \endfirsthead
  \toprule
  组件 & 组件功能 & AI使用情况 \\
  \midrule
  \endhead
  \code{wateros-task} & 进程、线程、调度与生命周期 &
  当前组件无AI生成内容。\\
  \code{wateros-base} & 基础类型、配置与同步原语 &
  当前组件无AI生成内容。\\
  \code{wateros-runtime} & 启动、日志、panic与内核堆支持 &
  当前组件无AI生成内容。\\
  \code{wateros-network} & IPv4、TCP/UDP与socket运行时 &
  当前组件无AI生成内容。\\
  \code{wateros-klog} & 内核日志记录与读取 &
  当前组件无AI生成内容。\\
  \code{wateros-utils} & 无状态通用工具 &
  当前组件无AI生成内容。\\
  \code{wateros-platform} & 架构、硬件与固件平台适配 &
  平台相关汇编代码以及LoongArch架构和QEMU LoongArch平台实现由AI生成，涉及启动入口、异常与
  中断入口、上下文相关汇编、架构寄存器操作和平台服务接线；其余平台接口与RISC-V实现不计入。\\
  \code{wateros-mm} & 地址空间与虚拟内存 &
  AI生成了部分用户地址空间、用户访问、页表、堆与mmap实现，并分别适配RISC-V Sv39和
  LoongArch64；相关代码通过双架构构建以及缺页、COW、mmap/munmap和地址空间回收场景验证。\\
  \code{wateros-syscall} & 系统调用与ABI适配 &
  AI生成了用户内存复制和错误转换等公共路径，并补充文件、网络、任务、内存、poll、SysV IPC、
  mount和ioctl等generic64 ABI实现；合入前按Linux参数、flag和errno语义检查，并通过LTP、
  BusyBox及定向用户程序验证。\\
  \code{wateros-vfs} & 统一文件访问接口 &
  AI生成了FS桥接、fd会话和页缓存中的主要实现，包括文件与目录句柄、fd注册表、cwd、文件锁、
  挂载表、符号链接、tmpfs和分页句柄；相关功能通过文件、目录、dup/fork、poll和写回路径验证。\\
  \code{wateros-fs} & 文件系统实现与挂载 &
  AI生成了ext4适配与自测、procfs路径解析和内容渲染、devfs设备节点管理以及rootfs挂载注册等
  功能；项目成员结合镜像读写、目录项、挂载和设备节点测试检查结果。\\
  \code{wateros-ipc} & 信号、futex、共享内存与等待队列 &
  实现层部分由AI生成，涉及futex等待与唤醒、信号状态、共享内存、pipe和waitqueue等对象的
  状态管理、任务阻塞及资源回收；IPC公共接口和系统调用ABI由项目成员核对并通过定向测试验证。\\
  \code{wateros-driver} & 设备驱动与块缓存 &
  AI生成了块缓存管理器以及RTC、null字符设备适配，项目成员检查设备注册、基本读写、缓存同步、
  边界返回值和平台构建。\\
  \code{wateros-cred} & 任务身份、凭证与权限状态 &
  实现层部分由AI生成，包括\code{impl-root}中的任务凭证侧表、UID/GID与附加组状态、fork复制、
  线程共享以及任务退出时的清理逻辑；公共接口和权限语义由项目成员审查确认。\\
  \code{wateros-tty} & 控制台与伪终端 &
  组件完整由AI生成，包括termios公共类型、控制台行规程、canonical/raw输入、回显与控制字符、
  阻塞读取、\code{VMIN/VTIME}、UNIX98 PTY、等待队列以及关闭和hangup处理。项目成员负责与
  VFS、信号、UART和系统调用层集成，并使用shell、PTY和Nano-X终端场景验证。\\
  \code{wateros-gui} & 窗口、控件、输入与软件合成 &
  组件完整由AI生成，包括GUI公共数据模型、窗口层级与焦点管理、控件交互、键鼠事件转换、
  软件光栅绘制、脏区合成及framebuffer刷新。项目成员负责显示和输入设备接入，并通过Nano-X
  与桌面演示场景验证。\\
  \code{wateros-debug} & 调试状态与故障诊断 &
  组件完整由AI生成，包括面向GDB的诊断ABI、每CPU状态快照、事件环、构建标识、锁持有关系
  跟踪及\code{TrackedMutex}。项目成员完成调试入口接入，并通过故障注入和多核场景检查诊断
  数据的一致性。\\
  组件README与源码注释 & 开发文档与代码说明 &
  项目组件目录中的README文档以及面向维护的源码注释由AI生成和整理，内容包括组件职责、调用
  方向、数据结构、并发约束、Linux ABI差距、故障定位和回归入口；项目成员依据当前源码和实际
  测试结果复核，文档描述不作为实现已经完成的唯一依据。\\
  \bottomrule
\end{longtable}
